\IfFileExists{stacks-project.cls}{%
\documentclass{stacks-project}
}{%
\documentclass{amsart}
}

% For dealing with references we use the comment environment
\usepackage{verbatim}
\newenvironment{reference}{\comment}{\endcomment}
%\newenvironment{reference}{}{}
\newenvironment{slogan}{\comment}{\endcomment}
\newenvironment{history}{\comment}{\endcomment}

% For commutative diagrams we use Xy-pic
\usepackage[all]{xy}

% We use 2cell for 2-commutative diagrams.
\xyoption{2cell}
\UseAllTwocells

% We use multicol for the list of chapters between chapters
\usepackage{multicol}

% This is generally recommended for better output
\usepackage{lmodern}
\usepackage[T1]{fontenc}

% For cross-file-references
\usepackage{xr-hyper}

% Package for hypertext links:
\usepackage{hyperref}

% For any local file, say "hello.tex" you want to link to please
% use \externaldocument[hello-]{hello}
\externaldocument[introduction-]{introduction}
\externaldocument[conventions-]{conventions}
\externaldocument[sets-]{sets}
\externaldocument[categories-]{categories}
\externaldocument[topology-]{topology}
\externaldocument[sheaves-]{sheaves}
\externaldocument[sites-]{sites}
\externaldocument[stacks-]{stacks}
\externaldocument[fields-]{fields}
\externaldocument[algebra-]{algebra}
\externaldocument[brauer-]{brauer}
\externaldocument[homology-]{homology}
\externaldocument[derived-]{derived}
\externaldocument[simplicial-]{simplicial}
\externaldocument[more-algebra-]{more-algebra}
\externaldocument[smoothing-]{smoothing}
\externaldocument[modules-]{modules}
\externaldocument[sites-modules-]{sites-modules}
\externaldocument[injectives-]{injectives}
\externaldocument[cohomology-]{cohomology}
\externaldocument[sites-cohomology-]{sites-cohomology}
\externaldocument[dga-]{dga}
\externaldocument[dpa-]{dpa}
\externaldocument[sdga-]{sdga}
\externaldocument[hypercovering-]{hypercovering}
\externaldocument[schemes-]{schemes}
\externaldocument[constructions-]{constructions}
\externaldocument[properties-]{properties}
\externaldocument[morphisms-]{morphisms}
\externaldocument[coherent-]{coherent}
\externaldocument[divisors-]{divisors}
\externaldocument[limits-]{limits}
\externaldocument[varieties-]{varieties}
\externaldocument[topologies-]{topologies}
\externaldocument[descent-]{descent}
\externaldocument[perfect-]{perfect}
\externaldocument[more-morphisms-]{more-morphisms}
\externaldocument[flat-]{flat}
\externaldocument[groupoids-]{groupoids}
\externaldocument[more-groupoids-]{more-groupoids}
\externaldocument[etale-]{etale}
\externaldocument[chow-]{chow}
\externaldocument[intersection-]{intersection}
\externaldocument[pic-]{pic}
\externaldocument[weil-]{weil}
\externaldocument[adequate-]{adequate}
\externaldocument[dualizing-]{dualizing}
\externaldocument[duality-]{duality}
\externaldocument[discriminant-]{discriminant}
\externaldocument[derham-]{derham}
\externaldocument[local-cohomology-]{local-cohomology}
\externaldocument[algebraization-]{algebraization}
\externaldocument[curves-]{curves}
\externaldocument[resolve-]{resolve}
\externaldocument[models-]{models}
\externaldocument[functors-]{functors}
\externaldocument[equiv-]{equiv}
\externaldocument[pione-]{pione}
\externaldocument[etale-cohomology-]{etale-cohomology}
\externaldocument[proetale-]{proetale}
\externaldocument[relative-cycles-]{relative-cycles}
\externaldocument[more-etale-]{more-etale}
\externaldocument[trace-]{trace}
\externaldocument[crystalline-]{crystalline}
\externaldocument[spaces-]{spaces}
\externaldocument[spaces-properties-]{spaces-properties}
\externaldocument[spaces-morphisms-]{spaces-morphisms}
\externaldocument[decent-spaces-]{decent-spaces}
\externaldocument[spaces-cohomology-]{spaces-cohomology}
\externaldocument[spaces-limits-]{spaces-limits}
\externaldocument[spaces-divisors-]{spaces-divisors}
\externaldocument[spaces-over-fields-]{spaces-over-fields}
\externaldocument[spaces-topologies-]{spaces-topologies}
\externaldocument[spaces-descent-]{spaces-descent}
\externaldocument[spaces-perfect-]{spaces-perfect}
\externaldocument[spaces-more-morphisms-]{spaces-more-morphisms}
\externaldocument[spaces-flat-]{spaces-flat}
\externaldocument[spaces-groupoids-]{spaces-groupoids}
\externaldocument[spaces-more-groupoids-]{spaces-more-groupoids}
\externaldocument[bootstrap-]{bootstrap}
\externaldocument[spaces-pushouts-]{spaces-pushouts}
\externaldocument[spaces-chow-]{spaces-chow}
\externaldocument[groupoids-quotients-]{groupoids-quotients}
\externaldocument[spaces-more-cohomology-]{spaces-more-cohomology}
\externaldocument[spaces-simplicial-]{spaces-simplicial}
\externaldocument[spaces-duality-]{spaces-duality}
\externaldocument[formal-spaces-]{formal-spaces}
\externaldocument[restricted-]{restricted}
\externaldocument[spaces-resolve-]{spaces-resolve}
\externaldocument[formal-defos-]{formal-defos}
\externaldocument[defos-]{defos}
\externaldocument[cotangent-]{cotangent}
\externaldocument[examples-defos-]{examples-defos}
\externaldocument[algebraic-]{algebraic}
\externaldocument[examples-stacks-]{examples-stacks}
\externaldocument[stacks-sheaves-]{stacks-sheaves}
\externaldocument[criteria-]{criteria}
\externaldocument[artin-]{artin}
\externaldocument[quot-]{quot}
\externaldocument[stacks-properties-]{stacks-properties}
\externaldocument[stacks-morphisms-]{stacks-morphisms}
\externaldocument[stacks-limits-]{stacks-limits}
\externaldocument[stacks-cohomology-]{stacks-cohomology}
\externaldocument[stacks-perfect-]{stacks-perfect}
\externaldocument[stacks-introduction-]{stacks-introduction}
\externaldocument[stacks-more-morphisms-]{stacks-more-morphisms}
\externaldocument[stacks-geometry-]{stacks-geometry}
\externaldocument[moduli-]{moduli}
\externaldocument[moduli-curves-]{moduli-curves}
\externaldocument[examples-]{examples}
\externaldocument[exercises-]{exercises}
\externaldocument[guide-]{guide}
\externaldocument[desirables-]{desirables}
\externaldocument[coding-]{coding}
\externaldocument[obsolete-]{obsolete}
\externaldocument[fdl-]{fdl}
\externaldocument[index-]{index}

% Theorem environments.
%
\theoremstyle{plain}
\newtheorem{theorem}[subsection]{Theorem}
\newtheorem{proposition}[subsection]{Proposition}
\newtheorem{lemma}[subsection]{Lemma}

\theoremstyle{definition}
\newtheorem{definition}[subsection]{Definition}
\newtheorem{example}[subsection]{Example}
\newtheorem{exercise}[subsection]{Exercise}
\newtheorem{situation}[subsection]{Situation}

\theoremstyle{remark}
\newtheorem{remark}[subsection]{Remark}
\newtheorem{remarks}[subsection]{Remarks}

\numberwithin{equation}{subsection}

% Macros
%
\def\lim{\mathop{\mathrm{lim}}\nolimits}
\def\colim{\mathop{\mathrm{colim}}\nolimits}
\def\Spec{\mathop{\mathrm{Spec}}}
\def\Hom{\mathop{\mathrm{Hom}}\nolimits}
\def\Ext{\mathop{\mathrm{Ext}}\nolimits}
\def\SheafHom{\mathop{\mathcal{H}\!\mathit{om}}\nolimits}
\def\SheafExt{\mathop{\mathcal{E}\!\mathit{xt}}\nolimits}
\def\Sch{\mathit{Sch}}
\def\Mor{\mathop{\mathrm{Mor}}\nolimits}
\def\Ob{\mathop{\mathrm{Ob}}\nolimits}
\def\Sh{\mathop{\mathit{Sh}}\nolimits}
\def\NL{\mathop{N\!L}\nolimits}
\def\CH{\mathop{\mathrm{CH}}\nolimits}
\def\proetale{{pro\text{-}\acute{e}tale}}
\def\etale{{\acute{e}tale}}
\def\QCoh{\mathit{QCoh}}
\def\Ker{\mathop{\mathrm{Ker}}}
\def\Im{\mathop{\mathrm{Im}}}
\def\Coker{\mathop{\mathrm{Coker}}}
\def\Coim{\mathop{\mathrm{Coim}}}

% Boxtimes
%
\DeclareMathSymbol{\boxtimes}{\mathbin}{AMSa}{"02}

%
% Macros for moduli stacks/spaces
%
\def\QCohstack{\mathcal{QC}\!\mathit{oh}}
\def\Cohstack{\mathcal{C}\!\mathit{oh}}
\def\Spacesstack{\mathcal{S}\!\mathit{paces}}
\def\Quotfunctor{\mathrm{Quot}}
\def\Hilbfunctor{\mathrm{Hilb}}
\def\Curvesstack{\mathcal{C}\!\mathit{urves}}
\def\Polarizedstack{\mathcal{P}\!\mathit{olarized}}
\def\Complexesstack{\mathcal{C}\!\mathit{omplexes}}
% \Pic is the operator that assigns to X its picard group, usage \Pic(X)
% \Picardstack_{X/B} denotes the Picard stack of X over B
% \Picardfunctor_{X/B} denotes the Picard functor of X over B
\def\Pic{\mathop{\mathrm{Pic}}\nolimits}
\def\Picardstack{\mathcal{P}\!\mathit{ic}}
\def\Picardfunctor{\mathrm{Pic}}
\def\Deformationcategory{\mathcal{D}\!\mathit{ef}}
\usepackage{cleveref}
\crefname{originallemma}{original lemma}{original lemmas}
\Crefname{originallemma}{Original lemma}{Original lemmas}
\makeatletter
\newlabel{lemma-directed-colimit-finite-type}{{0829}{source}{}{sourcepointer-1}{}}
\newlabel{lemma-directed-colimit-finite-type@cref}{{[originallemma][0][]0829}{[1][0][]source}{}{}{}}
\makeatother
\crefname{originalreference}{original reference}{original references}
\Crefname{originalreference}{Original reference}{Original references}
\makeatletter
\newlabel{properties-lemma-quasi-coherent-finite-type-ideals}{{01PH}{source}{}{sourcepointer-2}{}}
\newlabel{properties-lemma-quasi-coherent-finite-type-ideals@cref}{{[originalreference][0][]01PH}{[1][0][]source}{}{}{}}
\makeatother
\crefname{originalreference}{original reference}{original references}
\Crefname{originalreference}{Original reference}{Original references}
\makeatletter
\newlabel{spaces-morphisms-lemma-quasi-compact-local}{{03KG}{source}{}{sourcepointer-3}{}}
\newlabel{spaces-morphisms-lemma-quasi-compact-local@cref}{{[originalreference][0][]03KG}{[1][0][]source}{}{}{}}
\makeatother
\crefname{originalreference}{original reference}{original references}
\Crefname{originalreference}{Original reference}{Original references}
\makeatletter
\newlabel{spaces-morphisms-lemma-quasi-compact-quasi-separated-permanence}{{073B}{source}{}{sourcepointer-4}{}}
\newlabel{spaces-morphisms-lemma-quasi-compact-quasi-separated-permanence@cref}{{[originalreference][0][]073B}{[1][0][]source}{}{}{}}
\makeatother
\crefname{originalreference}{original reference}{original references}
\Crefname{originalreference}{Original reference}{Original references}
\makeatletter
\newlabel{spaces-properties-lemma-quasi-compact-affine-cover}{{03H6}{source}{}{sourcepointer-5}{}}
\newlabel{spaces-properties-lemma-quasi-compact-affine-cover@cref}{{[originalreference][0][]03H6}{[1][0][]source}{}{}{}}
\makeatother
\crefname{originalreference}{original reference}{original references}
\Crefname{originalreference}{Original reference}{Original references}
\makeatletter
\newlabel{spaces-properties-lemma-reduced-closed-subspace}{{03IQ}{source}{}{sourcepointer-6}{}}
\newlabel{spaces-properties-lemma-reduced-closed-subspace@cref}{{[originalreference][0][]03IQ}{[1][0][]source}{}{}{}}
\makeatother
\title{Quasi-compact opens of qcqs spaces cut out by finite type ideals}
\author{The Stacks Project Authors}
\date{stacks-project:spaces-limits.tex:local-snapshot-20261012; GNU FDL 1.2 or later}
\begin{document}
\maketitle
\noindent\textit{Editorial note.} Counted span is the author's own proof environment, source lines 3396-3431, following the statement at 3382-3394.\par\medskip
\noindent\textit{Editorial note.} No omissions; the surjectivity of |phi\textasciicircum{}\{-1\}U| -> |U| is justified in the text.\par\medskip
\noindent\textit{Allowed standard prerequisites (editorial).}\begin{itemize}
\item Affine etale cover of a qcqs space. Properties of Spaces, Lemma spaces-properties-lemma-quasi-compact-affine-cover (cross-chapter, cited at \textasciitilde{}3397-3399). Surjective etale phi : W -> X.
\item Quasi-compactness permanence for qcqs sources. Morphisms of Spaces, Lemma spaces-morphisms-lemma-quasi-compact-quasi-separated-permanence (cross-chapter, cited at \textasciitilde{}3403-3405). phi is quasi-compact since X is quasi-separated.
\item Scheme case of the ideal criterion. Properties, Lemma properties-lemma-quasi-coherent-finite-type-ideals (cross-chapter, cited at \textasciitilde{}3406-3407 and \textasciitilde{}3426-3429). Equivalence for the affine scheme W, used in both directions.
\item Quasi-compact is local on the target. Morphisms of Spaces, Lemma spaces-morphisms-lemma-quasi-compact-local (cross-chapter, cited at \textasciitilde{}3408-3409). Lifts quasi-compactness back to U -> X.
\item Reduced induced closed subspace ideal. Properties of Spaces, Lemma spaces-properties-lemma-reduced-closed-subspace (cross-chapter, cited at \textasciitilde{}3413-3417). The unique quasi-coherent ideal J cutting out the complement.
\item Approximating sections by finite type subsheaves. lemma-directed-colimit-finite-type (same chapter, cited at \textasciitilde{}3420-3424). Finite type quasi-coherent sub-ideal I agreeing with J on U.
\end{itemize}\medskip
\ifdefined\hypertarget\hypertarget{sourcepointer-1}{}\fi
\noindent\textit{Source reference.} Original lemma 0829: sources/stacks-spaces-limits/pinned-original-source.tex, source lines 2519--2519. Navigation pointer only: the counted proof cites this statement by \textbackslash{}ref as a standing input; the deduction the counted proof has to make is carried out inside the counted span, and no part of that deduction is delegated to this pointer.\par\medskip
\ifdefined\hypertarget\hypertarget{sourcepointer-2}{}\fi
\noindent\textit{Source reference.} Original reference 01PH: focused-seeds/geometry/S9-spaces-limits-quasi-coherent-finite-type-ideals/ref-tags.txt, source lines 5--5. Navigation pointer only: the counted proof cites this statement by \textbackslash{}ref as a standing input; the deduction the counted proof has to make is carried out inside the counted span, and no part of that deduction is delegated to this pointer.\par\medskip
\ifdefined\hypertarget\hypertarget{sourcepointer-3}{}\fi
\noindent\textit{Source reference.} Original reference 03KG: focused-seeds/geometry/S9-spaces-limits-quasi-coherent-finite-type-ideals/ref-tags.txt, source lines 6--6. Navigation pointer only: the counted proof cites this statement by \textbackslash{}ref as a standing input; the deduction the counted proof has to make is carried out inside the counted span, and no part of that deduction is delegated to this pointer.\par\medskip
\ifdefined\hypertarget\hypertarget{sourcepointer-4}{}\fi
\noindent\textit{Source reference.} Original reference 073B: focused-seeds/geometry/S9-spaces-limits-quasi-coherent-finite-type-ideals/ref-tags.txt, source lines 7--7. Navigation pointer only: the counted proof cites this statement by \textbackslash{}ref as a standing input; the deduction the counted proof has to make is carried out inside the counted span, and no part of that deduction is delegated to this pointer.\par\medskip
\ifdefined\hypertarget\hypertarget{sourcepointer-5}{}\fi
\noindent\textit{Source reference.} Original reference 03H6: focused-seeds/geometry/S9-spaces-limits-quasi-coherent-finite-type-ideals/ref-tags.txt, source lines 8--8. Navigation pointer only: the counted proof cites this statement by \textbackslash{}ref as a standing input; the deduction the counted proof has to make is carried out inside the counted span, and no part of that deduction is delegated to this pointer.\par\medskip
\ifdefined\hypertarget\hypertarget{sourcepointer-6}{}\fi
\noindent\textit{Source reference.} Original reference 03IQ: focused-seeds/geometry/S9-spaces-limits-quasi-coherent-finite-type-ideals/ref-tags.txt, source lines 9--9. Navigation pointer only: the counted proof cites this statement by \textbackslash{}ref as a standing input; the deduction the counted proof has to make is carried out inside the counted span, and no part of that deduction is delegated to this pointer.\par\medskip
\begin{lemma}
\label{lemma-quasi-coherent-finite-type-ideals}
Let $S$ be a scheme.
Let $X$ be a quasi-compact and quasi-separated algebraic space.
Let $U \subset X$ be an open subspace. The following are equivalent:
\begin{enumerate}
\item $U \to X$ is quasi-compact,
\item $U$ is quasi-compact, and
\item there exists a finite type quasi-coherent sheaf of ideals
$\mathcal{I} \subset \mathcal{O}_X$ such that
$|X| \setminus |U| = |V(\mathcal{I})|$.
\end{enumerate}
\end{lemma}

\begin{proof}
Let $W$ be an affine scheme and let $\varphi : W \to X$ be a surjective
\'etale morphism, see Properties of Spaces, Lemma
\ref{spaces-properties-lemma-quasi-compact-affine-cover}.
If (1) holds, then $\varphi^{-1}(U) \to W$ is quasi-compact, hence
$\varphi^{-1}(U)$ is quasi-compact, hence $U$ is quasi-compact
(as $|\varphi^{-1}(U)| \to |U|$ is surjective). If (2) holds, then
$\varphi^{-1}(U)$ is quasi-compact because $\varphi$ is quasi-compact
since $X$ is quasi-separated (Morphisms of Spaces,
Lemma \ref{spaces-morphisms-lemma-quasi-compact-quasi-separated-permanence}).
Hence $\varphi^{-1}(U) \to W$ is a quasi-compact morphism of schemes by
Properties, Lemma \ref{properties-lemma-quasi-coherent-finite-type-ideals}.
It follows that $U \to X$ is quasi-compact by
Morphisms of Spaces, Lemma \ref{spaces-morphisms-lemma-quasi-compact-local}.
Thus (1) and (2) are equivalent.

\medskip\noindent
Assume (1) and (2). By
Properties of Spaces, Lemma
\ref{spaces-properties-lemma-reduced-closed-subspace} there exists
a unique quasi-coherent sheaf of ideals $\mathcal{J}$ cutting
out the reduced induced closed subspace structure on $|X| \setminus |U|$.
Note that $\mathcal{J}|_U = \mathcal{O}_U$ which is an
$\mathcal{O}_U$-modules of finite type.
As $U$ is quasi-compact it follows from
Lemma \ref{lemma-directed-colimit-finite-type}
that there exists a quasi-coherent subsheaf
$\mathcal{I} \subset \mathcal{J}$ which is of finite type
and has the property that $\mathcal{I}|_U = \mathcal{J}|_U$.
Then $|X| \setminus |U| = |V(\mathcal{I})|$ and we obtain (3). Conversely,
if $\mathcal{I}$ is as in (3), then $\varphi^{-1}(U) \subset W$
is a quasi-compact open by the lemma for schemes
(Properties, Lemma \ref{properties-lemma-quasi-coherent-finite-type-ideals})
applied to $\varphi^{-1}\mathcal{I}$ on $W$.
Thus (2) holds.
\end{proof}

\end{document}
